두 이진 문자열 $A, B$를 비트별 배타적 논리합(bitwise XOR)한 값 $A \oplus B$는 다음과 같이 정의된다.
$A$의 $i$번째 자리 문자와 $B$의 $i$번째 자리 문자가 서로 다르면 $A \oplus B$의 $i$번째 자리 문자가 $1$이고, 서로 같으면 $A \oplus B$의 $i$번째 자리 문자는 $0$이다.
예를 들어 $A=\text{1100}$, $B=\text{1010}$이므로 $A \oplus B = \text{1100} \oplus \text{1010} = \text{0110}$으로 계산된다.
$k$개의 이진 문자열 $S_1, S_2, \cdots, S_k$를 비트별 배타적 논리합한 결과는 $(\dots((S_1 \oplus S_2) \oplus S_3) \oplus \dots) \oplus S_k$로 정의된다. 이 연산의 결과는 $S_1, S_2, \cdots, S_k$의 순서를 바꾸더라도 변하지 않음을 증명할 수 있다.

팰린드롬이란, 앞에서부터 읽으나 뒤에서부터 읽으나 같은 문자열을 말한다. 예를 들어, \texttt{010}, \texttt{0110} 등은 팰린드롬이며 \texttt{011}, \texttt{1011} 등은 팰린드롬이 아니다.